\begin{abstract}
Decoy tool calls change both an LLM Agent's encrypted traffic and the work its services perform. We study how call generation affects task recognition and measure inference, goal completion, business-state integrity and cost on the same executions. Our 3,000 single-host executions use two fixed-weight planners over MCP and TLS~1.3/HTTP/2. Controller-inserted cover bypasses the model wait preceding ordinary calls, exposing a short-gap signature. An exploratory reanalysis shows that time remains useful after waits are sampled from training traces: removing time features reduces selected-attack recognition in three of four history-dispatch cells after Holm correction, including 99 to 66 and 106 to 77 of 120 concurrent tasks. Context removal and joint decoding also recover recognition without supplied test call windows, while retaining episode association. They improve over refitted baselines in five of eight cells after Holm correction; concurrent scores rise from 56 to 80 and from 78 to 112/120. Neither shadow cell improves with the added attacks. Shadow cover instead has its own complete-policy trade-off: more and different calls, about twice the model tokens, and residual record-only maxima of 73 and 69/120. These results connect cover generation to observable trace structure and motivate evaluating its privacy, business authority and resource cost together.
\end{abstract}
